% Research fragment for the parent article. No standalone preamble.
% Source antecedent: incoming Nested Harmonic Jets, article.tex,
% proposition prop:zeroflow and Further research question 7.
% Prefix mz: avoids conflicts with the source manuscript labels.

\section{All spectral derivatives at the moving Gamma zeros}
\label{mz:sec:zeros}\label{sec:roots}

The incoming \emph{Nested Harmonic Jets} report \cite{nested} introduces a holomorphic
spectral product, evaluates its first spectral derivative at each zero of
its Gamma slice, and asks for an explicit construction at every derivative
order. We give such a construction. Its coefficients have a finite
triangular description by normalized Stieltjes antiderivatives. This solves
the algorithmic part of that question; it gives a sufficient coordinate
system rather than an arithmetic minimality theorem for constants.

The Weierstrass Gamma product, Bell polynomials, Stirling numbers, and
Mellin transforms used below are classical. The contribution here is their
specific combination into the all-order zero evaluation, its finite
Stieltjes primitive closure, and a harmonic-number coefficient algorithm.
The elementary digamma integral appearing at second order is not claimed
as a new evaluation in isolation.

Fix $a>0$, put $x_n=n+a$, and adopt
\[
 \zeta(1+u,a)=\frac1u+
       \sum_{r\ge0}\frac{(-1)^r\gamma_r(a)}{r!}u^r,
 \qquad \gamma_r=\gamma_r(1),\qquad \gamma_0(a)=-\psi(a).
\]
Use the source product and its specified counterterm:
\begin{align}
 \mathcal G(s,z;a)
 &=e^{z(\zeta(s,a)-1/(s-1))}
       \prod_{n\ge0}(1+zx_n^{-s})e^{-zx_n^{-s}},
       \qquad \Re s>\tfrac12,\label{mz:G}\\
 F_R(z;a)&=\left.\partial_s^R\mathcal G(s,z;a)\right|_{s=1}.
 \label{mz:F}
\end{align}
The pole-subtracted expression in \eqref{mz:G} is interpreted by its
holomorphic extension at one. Local uniform convergence, together with the
classical Gamma product, gives
\[
 \mathcal G(1,z;a)=\frac{\Gamma(a)}{\Gamma(a+z)}.
\]
We evaluate $F_R$ at the fixed argument $z=-a-m$. This is different from
differentiating along the zero path $z=-(a+m)^s$, where the function is
identically zero.

\subsection{A finite primitive coordinate system}

For $x\ge0$, integers $r\ge0$, and integers $k\ge2$, define
\begin{align}
 L_r(x)&=\sum_{n=1}^{\infty}
   \frac{\log^r(n+x)-\log^r n}{n},\qquad L_0(x)=0,
       \label{mz:L}\\
 M_{r,k}(x)&=\sum_{n=1}^{\infty}
          \frac{\log^r(n+x)}{n^k}.
       \label{mz:M}
\end{align}
The convention for the zeroth power is one. The sums and their fixed
derivatives converge locally uniformly: the summands of $L_r$, for $r>0$,
are $O(\log^{r-1}n/n^2)$ on a compact $x$-interval, and those of $M_{r,k}$
are $O(\log^r n/n^k)$. They also define holomorphic functions on
$\mathbb C\setminus(-\infty,-1]$ using principal logarithms.

Put
\begin{equation}
 T_r(x)=L_r(x)+\gamma_r-\gamma_r(1+x).
 \label{mz:T}
\end{equation}

\begin{theorem}[Finite Stieltjes primitive closure]
\label{mz:thm:primitive}
For $r\ge1$ and $x>0$,
\begin{align}
 L_r'(x)&=\frac r{x}T_{r-1}(x),\qquad L_r(0)=0,
      \label{mz:Lrec}\\
 M_{r,k}'(x)
 &=r\left\{\sum_{j=2}^{k}
     \frac{(-1)^{k-j}}{x^{k-j+1}}M_{r-1,j}(x)
       +\frac{(-1)^{k-1}}{x^k}T_{r-1}(x)\right\},
      \label{mz:Mrec}\\
 M_{r,k}(0)&=(-1)^r\zeta^{(r)}(k),
       \qquad M_{0,k}(x)=\zeta(k).
      \label{mz:Minitial}
\end{align}
The combined right sides have removable singularities at zero.
Consequently, any finite rectangle of $L_r$ and $M_{r,k}$ can be
constructed by finitely many ordinary integrations from zero, starting
with ordinary zeta derivatives and the generalized Stieltjes functions.
\end{theorem}

\begin{proof}
The defining limit for generalized Stieltjes constants gives
\begin{align}
 \sum_{n=1}^{\infty}\log^r(n+x)
          \left(\frac1n-\frac1{n+x}\right)
 &=L_r(x)+\gamma_r-\gamma_r(1+x)\notag\\
 &=T_r(x).
 \label{mz:Tsum}
\end{align}
To see the normalization explicitly, subtract the two cutoff sums for
$\gamma_r$ and $\gamma_r(1+x)$. Their integral counterterms differ by
$[\log^{r+1}(N+1+x)-\log^{r+1}(N+1)]/(r+1)$, which tends to zero.
The remaining difference is absolutely convergent. Differentiating
\eqref{mz:L} termwise and using
$1/[n(n+x)]=x^{-1}(1/n-1/(n+x))$ proves \eqref{mz:Lrec}.

For $k\ge2$, the elementary partial-fraction identity
\begin{equation}
 \frac1{n^k(n+x)}
 =\sum_{j=2}^{k}\frac{(-1)^{k-j}}{x^{k-j+1}n^j}
   +\frac{(-1)^{k-1}}{x^k}
          \left(\frac1n-\frac1{n+x}\right)
 \label{mz:partial}
\end{equation}
and termwise differentiation of \eqref{mz:M} prove \eqref{mz:Mrec}.
At zero, differentiation of the absolutely convergent Dirichlet series
for $\zeta(k)$ gives \eqref{mz:Minitial}. The original differentiated
series are holomorphic near zero. Thus the combined rational expressions
on the right of \eqref{mz:Lrec}--\eqref{mz:Mrec} have removable
singularities there. Integrating those combined expressions proves the
last assertion by induction on $r$.
\end{proof}

One must integrate the combined right side of \eqref{mz:Mrec}. Its
individual displayed terms can have nonintegrable poles at zero, although
their sum is regular.

In particular, the first coordinate is the ordinary digamma primitive
\begin{equation}
 D(x):=L_1(x)
 =\sum_{n\ge1}\frac{\log(1+x/n)}n
 =\int_0^x\frac{\psi(1+t)+\gamma}{t}\,dt.
 \label{mz:D}
\end{equation}
The integrand takes the limiting value $\zeta(2)$ at zero. Higher $L_r$
are obtained from
\[
 L_r(x)=r\int_0^x
       \frac{L_{r-1}(t)+\gamma_{r-1}-\gamma_{r-1}(1+t)}t\,dt.
\]
For example
\begin{align}
 L_2'(x)&=\frac2x\bigl(D(x)+\gamma_1-\gamma_1(1+x)\bigr),\\
 M_{2,2}'(x)
 &=\frac2xM_{1,2}(x)
   -\frac2{x^2}\bigl(D(x)+\gamma_1-\gamma_1(1+x)\bigr).
\end{align}

\subsection{The zero-factor theorem}

Let $\left\{\begin{smallmatrix}r\\k\end{smallmatrix}\right\}$ be a
Stirling number of the second kind and set
\begin{equation}
 A_{r,k}=(k-1)!
        \left\{\begin{matrix}r\\k\end{matrix}\right\}.
 \label{mz:A}
\end{equation}
Let $B_j$ be the Bernoulli numbers with $B_1=-1/2$. Define the complete
exponential Bell polynomials by
\[
 \exp\left(\sum_{j\ge1}c_j\frac{u^j}{j!}\right)
       =\sum_{n\ge0}Y_n(c_1,\ldots,c_n)\frac{u^n}{n!}.
\]

\begin{theorem}[All fixed-argument spectral jets at a Gamma zero]
\label{mz:thm:zero}
Let $a>0$, $m\in\mathbb Z_{\ge0}$, and set
\[
 x=a+m,\qquad \ell=\log x,\qquad
 h_0=x\Gamma(a)(-1)^m m!.
\]
For $r\ge1$ define the finite expression
\begin{align}
 b_r={}&(-1)^{r+1}\left\{
      x\bigl(\gamma_r+L_r(x)\bigr)
       +\sum_{k=2}^{r}A_{r,k}x^kM_{r,k}(x)\right.\notag\\
 &\hspace{25mm}\left.
       +\sum_{j=1}^{m}\log^r(x-j)
            \sum_{k=1}^{r}A_{r,k}\left(-\frac xj\right)^k
                        \right\},\label{mz:b}\\
 c_r={}&b_r+\frac{B_r}{r}\ell^r.
 \label{mz:c}
\end{align}
Empty sums vanish. Then, for every $R\ge1$,
\begin{equation}
 \boxed{F_R(-x;a)
    =R h_0\ell\,Y_{R-1}(c_1,\ldots,c_{R-1}).}
 \label{mz:allzero}
\end{equation}
At $x=1$, every $F_R(-1;1)$ is zero. For $x\ne1$, the formula requires
only the $R-1$ logarithmic jets $b_1,\ldots,b_{R-1}$.
\end{theorem}

\begin{proof}
Remove the unique vanishing factor:
\[
 \mathcal G(1+u,-x;a)=(1-e^{-\ell u})\mathcal H(1+u),
\]
where $\mathcal H$ is holomorphic and nonzero near one. At $s=1$, the
Gamma slice has derivative $\Gamma(a)(-1)^m m!$ in $z$ at $-x$.
The removed factor is $(z+x)/x$ on that slice, so $\mathcal H(1)=h_0$.
Equivalently, for a cutoff $N>m$, the zero-removed value is
\begin{equation}
 \mathcal H_N(1)=h_0 (N+a)^x
                      \frac{\Gamma(N-m)}{\Gamma(N+a)},
 \label{mz:Hcut}
\end{equation}
which tends to $h_0$ by the Gamma ratio asymptotic.

The source cutoff is
\[
 \mathcal G_N(s,-x;a)
   =e^{xC_N(s;a)}\prod_{n<N}(1-xx_n^{-s}),\qquad
 C_N(s;a)=\frac{1-(N+a)^{1-s}}{s-1}.
\]
Its derivatives satisfy
$\partial_s^r C_N(1;a)=(-1)^r\log^{r+1}(N+a)/(r+1)$.
For integer $r\ge1$,
\begin{equation}
 \left.\partial_s^r\log(1-xx_n^{-s})\right|_{s=1}
       =(-1)^{r+1}\log^r x_n\,
                         \operatorname{Li}_{1-r}(x/x_n).
 \label{mz:LiDerivative}
\end{equation}
Only integer nonpositive polylogarithm orders occur here, hence these are
rational functions; the formula is valid even for $n<m$, when the
argument exceeds one. The rational identity
\begin{equation}
 \operatorname{Li}_{1-r}(q)
       =\sum_{k=1}^{r}A_{r,k}\left(\frac q{1-q}\right)^k
 \label{mz:LiStirling}
\end{equation}
follows either by applying $q\partial_q$ repeatedly to $q/(1-q)$ or
from the Stirling recurrence. Thus
\begin{align}
 \left.\partial_s^r\log\mathcal H_N(s)\right|_{s=1}
 =(-1)^r\left\{\frac{x}{r+1}\log^{r+1}(N+a)
  -\sum_{\substack{0\le n<N\\n\ne m}}\log^r(n+a)
          \sum_{k=1}^{r}\frac{A_{r,k}x^k}{(n-m)^k}\right\}.
 \label{mz:bcut}
\end{align}
The $n<m$ terms become the finite $j$-sum in \eqref{mz:b}. Set $n=m+j$
in the remaining terms. Since $A_{r,1}=1$,
\[
 \lim_{J\to\infty}\left\{
   \sum_{j=1}^{J}\frac{\log^r(j+x)}j
       -\frac{\log^{r+1}(J+1+x)}{r+1}\right\}
       =\gamma_r+L_r(x).
\]
Indeed, subtract the analogous limit for $\gamma_r$; the change in its
logarithmic counterterm tends to zero. The sums for $k\ge2$ converge to
$M_{r,k}(x)$. Local uniform convergence of the zero-removed products
justifies passage to their logarithmic derivatives, proving
$b_r=\partial_s^r\log\mathcal H(1)$ and \eqref{mz:b}.

For $\ell\ne0$, the Bernoulli generating function yields, near $u=0$,
\begin{equation}
 \log\frac{1-e^{-\ell u}}{\ell u}
       =\sum_{r\ge1}\frac{B_r\ell^r}{r}\frac{u^r}{r!}.
 \label{mz:Bernoulli}
\end{equation}
Consequently
\[
 \mathcal G(1+u,-x;a)
       =h_0\ell u\exp\left(\sum_{r\ge1}c_r\frac{u^r}{r!}\right).
\]
Bell coefficient extraction proves \eqref{mz:allzero}. If $\ell=0$,
the removed factor $1-e^{-\ell u}$ vanishes identically, proving the
stated exceptional case without division by zero.
\end{proof}

The first three nontrivial rows are
\begin{align}
 F_1(-x;a)&=h_0\ell,\label{mz:F1}\\
 F_2(-x;a)&=h_0(2\ell b_1-\ell^2),\label{mz:F2b}\\
 F_3(-x;a)&=h_0\{3\ell(b_1^2+b_2)-3\ell^2b_1+\ell^3\}.
 \label{mz:F3}
\end{align}
The first is the source's already proved zero-flow formula. The second has
the explicit Gamma--Stieltjes--digamma-primitive evaluation
\begin{equation}
 \boxed{F_2(-x;a)=h_0\left[
   2x\ell\left\{\gamma_1+D(x)
           -\sum_{j=1}^{m}\frac{\log(x-j)}j\right\}-\ell^2\right].}
 \label{mz:F2}
\end{equation}
For the next row,
\begin{align}
 b_2=-\biggl[&x(\gamma_2+L_2(x))+x^2M_{2,2}(x)\notag\\
 &+\sum_{j=1}^{m}\log^2(x-j)
                    \left(-\frac xj+\frac{x^2}{j^2}\right)\biggr].
 \label{mz:b2}
\end{align}

For fixed $R\ge2$, Theorem~\ref{mz:thm:primitive} supplies a finite construction
of all entries needed in \eqref{mz:allzero}. It uses ordinary Stieltjes
constants only through $\gamma_{R-1}$, Stieltjes \emph{functions}
$\gamma_j(1+t)$ only through $j=R-2$, and finitely many zeta derivatives
$\zeta^{(r)}(k)$ with $0\le r\le R-1$ and $2\le k\le R-1$.
For $R\ge2$, the highest ordinary Stieltjes constant occurs linearly,
with coefficient
\begin{equation}
 [\gamma_{R-1}]F_R(-x;a)
       =(-1)^R R x h_0\log x
 \label{mz:highest}
\end{equation}
in this explicit coordinate formula. This is a statement about the
displayed construction. It does not assert arithmetic independence of
Stieltjes constants from zeta derivatives or their integrals.

\subsection{Polylogarithmic Mellin realization}

The coordinate system also has a direct polylogarithm interpretation. Set
\begin{equation}
 J_k(v,x)=\sum_{n=1}^{\infty}\frac1{n^k(n+x)^v}.
 \label{mz:J}
\end{equation}
For $x\ge0$ and $\Re v>0$, termwise integration of the absolutely
convergent polylogarithm series gives
\begin{equation}
 J_k(v,x)=\frac1{\Gamma(v)}\int_0^{\infty}
     t^{v-1}e^{-xt}\operatorname{Li}_k(e^{-t})\,dt,
       \qquad k\ge1.
 \label{mz:Mellin}
\end{equation}
For $k\ge2$ its Dirichlet series is holomorphic near $v=0$, and
\[
 M_{r,k}(x)=(-1)^r\partial_v^rJ_k(0,x).
\]
For $k=1$, the difference
\[
 J_1(v,x)-\zeta(1+v)
 =\sum_{n\ge1}\frac{(n+x)^{-v}-n^{-v}}n
\]
is holomorphic near zero because its summands are $O(n^{-2-\Re v})$
there, uniformly on compact sets. It follows that
\begin{equation}
 J_1(v,x)=\frac1v+
       \sum_{r\ge0}\frac{(-1)^r(\gamma_r+L_r(x))}{r!}v^r.
 \label{mz:JLaurent}
\end{equation}
Thus the exact Gamma-zero derivatives are finite combinations of resonant
Laurent coefficients and regular spectral derivatives of the elementary
polylogarithmic Mellin transform \eqref{mz:Mellin}. This is a specific
shifted Dirichlet series, not an identification of its Laurent coefficients
with the ordinary Hurwitz Stieltjes constants.

\subsection{Generalized harmonic numbers give every coefficient}

Let $e_j(h)$ be defined by
\begin{equation}
 \prod_{n=1}^{h}(1+u/n)=\sum_{j=0}^{h}e_j(h)u^j,
 \qquad e_0(h)=1.
 \label{mz:e}
\end{equation}
These are polynomials in the generalized harmonic numbers
$H_h^{(j)}=\sum_{n=1}^h n^{-j}$:
\[
 e_1(h)=H_h,\qquad
 e_2(h)=\tfrac12\bigl(H_h^2-H_h^{(2)}\bigr),\qquad
 \sum_j e_j(h)u^j
  =\exp\left(\sum_{j\ge1}\frac{(-1)^{j-1}H_h^{(j)}}j u^j\right).
\]
For $r\ge1$, $h\ge1$, and $k\ge1$, put
\begin{equation}
 Q_{r,k,h}(q)
   =\sum_{j=1}^{\min(r,h)}\frac{r!}{(r-j)!}
           e_{j-1}(h-1)\zeta^{(r-j)}(k+h,q).
 \label{mz:Q}
\end{equation}

\begin{theorem}[Harmonic--Hurwitz coefficient formulas]
\label{mz:thm:coefficients}
Let $r\ge1$, let $N\ge0$ be an integer, and let
$x\in\mathbb C\setminus(-\infty,-1]$ satisfy $|x|<N+1$. Then
\begin{align}
 L_r(x)={}&\sum_{n=1}^{N}
        \frac{\log^r(n+x)-\log^r n}{n}
   +(-1)^r\sum_{h\ge1}\frac{(-x)^h}{h}Q_{r,1,h}(N+1),
       \label{mz:Lseries}\\
 M_{r,k}(x)={}&\sum_{n=1}^{N}\frac{\log^r(n+x)}{n^k}
       +(-1)^r\zeta^{(r)}(k,N+1)\notag\\
 &+(-1)^r\sum_{h\ge1}\frac{(-x)^h}{h}Q_{r,k,h}(N+1),
       \qquad k\ge2.
       \label{mz:Mseries}
\end{align}
The series converge normally on compact subsets of $|x|<N+1$.
Every positive $x$ is covered after a finite cutoff. The only harmonic
orders needed for a spectral derivative of order $r$ are at most $r-1$.
\end{theorem}

\begin{proof}
Split \eqref{mz:J} at $N$ and use the binomial expansion in the tail:
\[
 \sum_{n>N}\frac1{n^k(n+x)^v}
   =\sum_{h\ge0}\frac{(-1)^h(v)_h x^h}{h!}
                                      \zeta(k+v+h,N+1).
\]
For $|x|<N+1$, the double series and its fixed $v$ derivatives converge
normally on the appropriate local domain; the $h=0,k=1$ pole is treated
separately. Since
\[
 \frac{(v)_h}{h!}=\frac vh
       \prod_{j=1}^{h-1}(1+v/j)
       =\frac1h\sum_{j=1}^{h}e_{j-1}(h-1)v^j,
\]
Leibniz differentiation at zero gives the coefficient
$Q_{r,k,h}(N+1)/h$. Formula \eqref{mz:JLaurent} or the derivative
description of $M_{r,k}$ then proves \eqref{mz:Lseries}--\eqref{mz:Mseries}.
The finite harmonic expression follows from \eqref{mz:e}.
\end{proof}

For example, for $|x|<1$,
\begin{align}
 D(x)&=\sum_{h\ge1}\frac{(-1)^{h+1}x^h}{h}\zeta(h+1),\\
 L_2(x)&=2\sum_{h\ge1}\frac{(-x)^h}{h}
       \bigl\{\zeta'(h+1)+H_{h-1}\zeta(h+1)\bigr\},\\
 M_{2,2}(x)&=\zeta''(2)+2\sum_{h\ge1}\frac{(-x)^h}{h}
       \bigl\{\zeta'(h+2)+H_{h-1}\zeta(h+2)\bigr\}.
\end{align}

\paragraph{Precise remaining questions.}
The finite coordinate construction is complete, but whether some of its
nontrivial fixed-argument integrals reduce to a smaller field of special
values remains open. The first test is $D(x)$ at a positive integer or a
nonintegral rational point. No failure of a symbolic or numerical search
would prove non-reducibility. A second target is a global description of
the individual higher spectral jets as entire functions of $z$, with these
Gamma-zero evaluations as interpolation data. A third target is the
compatible construction for several simultaneously moving zero factors
in multicolour spectral products.
